% from ega1-10-fr.tex
\oldpage[I]{188}

\begin{proposition}[10.5.6]
\label{I.10.5.6-fr}
Soient $\mathfrak{X}$, $\mathfrak{Y}$ deux préschémas formels,
$\mathscr{J}$ (resp. $\mathscr{K}$) un Idéal de définition de
$\mathfrak{X}$ (resp. $\mathfrak{Y}$),
$f:\mathfrak{X}\to\mathfrak{Y}$ un morphisme de préschémas formels.
\begin{enumerate}
  \item[\rm{(i)}] Si
    $f^*(\mathscr{K})\mathscr{O}_{\mathfrak{X}}\subset\mathscr{J}$,
    il existe un unique morphisme
    $f':(\mathfrak{X},\mathscr{O}_{\mathfrak{X}}/\mathscr{J})\to
    (\mathfrak{Y},\mathscr{O}_{\mathfrak{Y}}/\mathscr{K})$ de
    préschémas usuels rendant commutatif le diagramme
    \[
      \label{I.10.5.6.1-fr}
      \xymatrix{
        (\mathfrak{X},\mathscr{O}_{\mathfrak{X}})\ar[r]^f &
        (\mathfrak{Y},\mathscr{O}_{\mathfrak{Y}})\\
        (\mathfrak{X},\mathscr{O}_{\mathfrak{X}}/\mathscr{J})
          \ar[r]_{f'}\ar[u] &
        (\mathfrak{Y},\mathscr{O}_{\mathfrak{Y}}/\mathscr{K})\ar[u]
      }
      \tag{10.5.6.1}
    \]
    où les flèches verticales sont les morphismes canoniques.
  \item[\rm{(ii)}] Supposons que
    $\mathfrak{X}=\operatorname{Spf}(A)$,
    $\mathfrak{Y}=\operatorname{Spf}(B)$ soient des schémas formels
    affines, $\mathscr{J}=\mathfrak{J}^{\Delta}$,
    $\mathscr{K}=\mathfrak{K}^{\Delta}$, où $\mathfrak{J}$ (resp.
    $\mathfrak{K}$) est un idéal de définition de $A$ (resp. $B$), et
    $f=({}^a\varphi,\widetilde{\varphi})$, où
    $\varphi:B\to A$ est un homomorphisme continu~; pour que
    $f^*(\mathscr{K})\mathscr{O}_{\mathfrak{X}}\subset\mathscr{J}$,
    il faut et il suffit que
    $\varphi(\mathfrak{K})\subset\mathfrak{J}$, et $f'$ est alors le
    morphisme $({}^a\varphi',\widetilde{\varphi'})$, où
    $\varphi':B/\mathfrak{K}\to A/\mathfrak{J}$ est l'homomorphisme
    déduit de $\varphi$ par passage aux quotients.
\end{enumerate}
\end{proposition}

\begin{proof}
\medskip\noindent
\begin{enumerate}
  \item[\rm{(i)}] Si $f=(\psi,\theta)$, l'hypothèse entraîne que
    l'image par
    $\theta^\sharp:\psi^*(\mathscr{O}_{\mathfrak{Y}})\to
    \mathscr{O}_{\mathfrak{X}}$ du faisceau d'idéaux
    $\psi^*(\mathscr{K})$ de $\psi^*(\mathscr{O}_{\mathfrak{Y}})$ est
    contenue dans $\mathscr{J}$
    (\hyperref[0.4.3.5-fr]{0, 4.3.5}). Par passage aux quotients, on
    déduit donc de $\theta^\sharp$ un homomorphisme de faisceau d'anneaux
    \[
      \omega:\psi^*(\mathscr{O}_{\mathfrak{Y}}/\mathscr{K})
      =\psi^*(\mathscr{O}_{\mathfrak{Y}})/\psi^*(\mathscr{K})
      \to\mathscr{O}_{\mathfrak{X}}/\mathscr{J}~;
    \]
    en outre, comme pour tout $x\in\mathfrak{X}$,
    $\theta_x^\sharp$ est un homomorphisme \emph{local}, il en est de
    même de $\omega_x$. Le morphisme d'espaces annelés
    $(\psi,\omega^b)$ est donc
    (\hyperref[I.2.2.1-fr]{2.2.1}) l'unique morphisme $f'$ d'espaces
    annelés répondant à la question.
  \item[\rm{(ii)}] La correspondance canonique fonctorielle entre
    morphismes de préschémas formels affines et homomorphismes continus
    d'anneaux (\hyperref[I.10.2.2-fr]{10.2.2}) montre que dans le cas
    considéré, la relation
    $f^*(\mathscr{K})\mathscr{O}_{\mathfrak{X}}\subset\mathscr{J}$
    entraîne que l'on a
    $f'=({}^a\varphi',\widetilde{\varphi'})$, où
    $\varphi':B/\mathfrak{K}\to A/\mathfrak{J}$ est l'unique
    homomorphisme rendant commutatif le diagramme
    \[
      \label{I.10.5.6.2-fr}
      \xymatrix{
        B\ar[r]^\varphi\ar[d] & A\ar[d]\\
        B/\mathfrak{K}\ar[r]_{\varphi'} & A/\mathfrak{J}
      }
      \tag{10.5.6.2}
    \]
    L'existence de $\varphi'$ implique donc que
    $\varphi(\mathfrak{K})\subset\mathfrak{J}$. Inversement, si cette
    condition est vérifiée, en désignant par $\varphi'$ l'unique
    homomorphisme rendant commutatif le diagramme
    (\hyperref[I.10.5.6.2-fr]{10.5.6.2}) et posant
    $f'=({}^a\varphi',\widetilde{\varphi'})$, il est clair que le
    diagramme (\hyperref[I.10.5.6.1-fr]{10.5.6.1}) est commutatif~; la
    considération des homomorphismes
    ${}^a\varphi^*(\mathscr{O}_{\mathfrak{Y}})\to
    \mathscr{O}_{\mathfrak{X}}$ et
    ${}^a{\varphi'}^*(\mathscr{O}_{\mathfrak{Y}}/\mathscr{K})\to
    \mathscr{O}_{\mathfrak{X}}/\mathscr{J}$ correspondant à $f$ et
    $f'$ respectivement, montre alors que cela entraîne la relation
    $f^*(\mathscr{K})\mathscr{O}_{\mathfrak{X}}\subset\mathscr{J}$.
\end{enumerate}
\end{proof}

Il est clair que la correspondance $f\mapsto f'$ définie ci-dessus est
\emph{fonctorielle}.

\subsection{Préschémas formels comme limites inductives de préschémas.}
\label{subsection:I.10.6-fr}

\begin{env}[10.6.1]
\label{I.10.6.1-fr}
Soient $\mathfrak{X}$ un préschéma formel,
$(\mathscr{J}_\lambda)$ un système fondamental d'Idéaux de définition de
$\mathfrak{X}$~; pour tout $\lambda$, soit $f_\lambda$ le morphisme
canonique
$(\mathfrak{X},\mathscr{O}_{\mathfrak{X}}/\mathscr{J}_\lambda)
\to\mathfrak{X}$ (\hyperref[I.10.5.2-fr]{10.5.2})~; pour
$\mathscr{J}_\mu\subset\mathscr{J}_\lambda$, l'homomorphisme canonique
$\mathscr{O}_{\mathfrak{X}}/\mathscr{J}_\mu\to
\mathscr{O}_{\mathfrak{X}}/\mathscr{J}_\lambda$ définit un morphisme

